\par \Statement Пусть $G$ — конечная группа порядка $2^n, n \in \N$. Тогда в $G$ есть подгруппа $H$ такая, что $[G:H]=2$.
\par \Proof Проведем индукцию по $n$. База, $n=1$, тривиальна, докажем переход.
\par Группа $G$ является $p$-группой, поэтому ее центр $Z(G)$ нетривиален. Поскольку $Z(G) \trianglelefteq G$,
можно рассмотреть группу $G/Z(G)$. Это 2-группа меньшего порядка, и по предположению индукции в $G/Z(G)$ есть подгруппа $H$ такая, что $[G/Z(G) : H]=2$. По теореме о
соответствии, $H=K/Z(G)$ для некоторой группы $K$ такой, что $Z(G) \leq K \leq G$. Значит, $[G:K]=2$, и переход доказан.

\par \Cor Пусть $G$ — группа, $|G|=2^n$. Тогда в $G$ существует такая цепочка подгрупп $G=G_n \geq G_{n-1} \geq \ldots \geq G_0 = \{e\}$, что $\forall i \in \{0, \ldots, n-1\}: [G_{i+1}:G_i]=2$.

\par \Th Точку $z \in \Comp$ можно построить циркулем и линейкой $\Leftrightarrow$ поле разложения ее минимального многочлена $m_z \in \Q[x]$ имеет степень $2^n, n \in \N$.
\par \textbf{Доказательство достаточности (справа налево):} \Proof Обозначим через $L$ поле разложения многочлена $m_z$.
Пусть $[L : \Q]=2^n, n \in \N$. Тогда в группе $G:=Aut_\Q L$ порядка $2^n$ существует такая цепочка подгрупп $G=G_n \geq \ldots \geq G_0=\{e\}$, что для всех $i \in \{0, \ldots, n-1\}$ выполнено $[G_{i+1} : G_i]=2$. Выберем соответствующую цепочке подгрупп башню расширений $\Q=L_n \subset \ldots \subset L_0=L$, тогда $\forall i \in \{0, \ldots, n-1\}: [L_i : L_{i+1}]=2$. Но $x \in L$, поэтому $x$ можно построить по уже доказанному критерию (см. билет 7 на отл). \EndProof
\par TODO возможно нужно доказательство в другую сторону?